% ============================================================================
%  SECTION 8 — Hands-on exercises
% ============================================================================
\begin{frame}[t]{}
  \vfill\begin{center}{\Huge\color{YaleBlue} 8.\ \ Explore it yourself}\\[0.5em]
  {\color{BoxGray} small exercises --- one per theory}\end{center}\vfill
\end{frame}

\begin{frame}[t]{Small exercises: reproduce the outcomes yourself}
  Every result in this lecture is generated by the \texttt{gr} package. The
  exercises expose a few tunables (``\texttt{YOUR TURN}'' blocks) so you can drive
  the same tissue and watch each theory respond.
  \vspace{0.3em}
  \small
  \begin{center}
  \renewcommand{\arraystretch}{1.35}
  \begin{tabular}{@{}p{0.20\textwidth}p{0.30\textwidth}p{0.42\textwidth}@{}}
    \toprule
    \textbf{Exercise} & \textbf{Change\ldots} & \textbf{\ldots and observe} \\
    \midrule
    \texttt{ex02} kinematic & growth gain $k_g$, pressure step, geometry &
      stress returns to the \emph{prescribed} set-point; crank pressure on the
      artery to trigger runaway --- the bar never does \\
    \texttt{ex03} full CMM & turnover time $1/k_d$, gain $K_\sigma$; hypertension
      vs.\ aneurysm & how fast the wall adapts; collagen/SMC take over the load as
      elastin is degraded \\
    \texttt{ex04} homogenized & scenario (hypertension / aneurysm) &
      the homogenized curve overlays the full CMM --- at a fraction of the runtime \\
    \texttt{ex05} equilibrated & surviving-elastin fraction $s$ &
      transient converges to the equilibrated $\lambda^*$; find the \emph{critical}
      elastin loss where the root disappears \\
    \texttt{ex06} stability & gain, collagen stiffness, deposition stretch $G$ &
      reproduce the stability map, then watch the ``aneurysm'' region shrink \\
    \bottomrule
  \end{tabular}
  \end{center}
  \srccite{Exercises in the \texttt{gr} repo: exercises/ex02--ex06; run with \texttt{uv run python exercises/ex0N\_*.py}}
\end{frame}

\begin{frame}[t]{The four learnings, as things to verify}
  \begin{enumerate}
    \item \textbf{Kinematic growth} (\texttt{ex02}): it either reaches the
    \emph{a-priori prescribed} homeostasis, or is unstable --- stability is
    imposed. \emph{Verify:} the stress plateau always lands on the set-point;
    push the artery to make it diverge.
    \item \textbf{Constrained mixture} (\texttt{ex03}): stability
    \emph{depends on the insult}. \emph{Verify:} mild hypertension adapts; a
    severe elastin insult runs away.
    \item \textbf{Homogenized} (\texttt{ex04}): a \emph{similar time trajectory}
    to the full model. \emph{Verify:} overlay the two and measure the gap ---
    and the speed-up.
    \item \textbf{Equilibrated} (\texttt{ex05}, \texttt{ex06}): \emph{matches}
    the transient theories \emph{if} an equilibrium exists (and does \emph{not}
    when it does not). \emph{Verify:} sweep $s$ to the critical loss where the
    equilibrated root vanishes, then confirm the transient diverges just past it.
  \end{enumerate}
\end{frame}

